\section{Introduction}\label{sec:introduction}

G\"odel's incompleteness theorems show that no sufficiently expressive effective
formal system closes over all arithmetic truth, and that no such system proves
its own global consistency in the usual way~\cite{Godel1931,HilbertBernays1939}. What they do not by
themselves provide is a structural language for theory growth: what is held
fixed, what changes, and under which ledger distinct changes are compared.
Foundational discussion often shifts among axiom strength, maximality, and
informal mathematical practice without fixing that structural
vocabulary~\cite{Maddy1988,SteelMaximizeUnify}. This paper proposes a more explicit
alternative.

The paper treats a formal theory as a packaged closure object rather than as a
bare list of axioms. Once that shift is made, a basic distinction appears.
Iterating a fixed idempotent package is not the same operation as changing the
package. The first saturates; the second is the only source of persistent strict
growth. This distinction is structural before it is arithmetical.

Across the Six-Birds line, six primitive operations recur: rewrite (\(P1\)),
gating or support restriction (\(P2\)), protocol or route mismatch (\(P3\)),
staging/indexing/sectors (\(P4\)), packaging/closure (\(P5\)), and
accounting/audit (\(P6\)). The present paper uses them asymmetrically. \(P5\) is
the fixed-package object. \(P1/P2\) are the mechanisms of package change. \(P6\)
fixes the frozen evaluation regime in which yield, cost, and efficiency are
assessed. \(P4\) supplies the support/index axis on which representative
comparison is meaningful. \(P3\) remains diagnostic-only: relevant for route
mismatch, but not theorem-bearing in the main argument.

The body now proceeds in theorem-paper order. Section~\ref{sec:primitive-roles}
gives explicit definitions of closure packages, persistent growth witnesses,
frozen slices, support cones, dominance, efficiency, cone agreement, slice
invariance, and representative alignment targets. Section~\ref{sec:fixed-package}
states and proves the base and bridge theorems: fixed-package saturation and
package-change necessity. Section~\ref{sec:frozen-frontier} states the
comparative results for frozen-slice dominance, efficiency transfer, and
restricted alignment. Section~\ref{sec:conditional-lift} states the main
conditional arithmetic theorem under an explicit assumption box and external
contract.

The resulting manuscript is theorem-led, but conditionally so. Its proof-carrying
core is internal and Lean-backed~\cite{deMouraUllrich2021,Mathlib2020}. Its arithmetic lift is conditional on an
explicit, finite external dependency contract. The vendor-backed PICA evidence
plays a different role: it supports the primitive-role assignment and the scope
limits of the manuscript, but it does not cross the proof boundary.

\paragraph{Contributions.}
The paper makes four contributions. First, it gives explicit body-level
mathematical definitions for the packaged closure objects, frozen-slice
comparison regime, and representative-alignment target used in the theorem route.
Second, it proves an internal theorem ladder: fixed-package saturation,
package-change necessity, and restricted alignment on a frozen slice. Third, it
states and derives a conditional arithmetic lift under an explicit external
dependency contract, thereby exposing rather than hiding the remaining arithmetic
boundary. Fourth, it uses vendor-backed PICA evidence only to support
primitive-role assignment and scope limits, not as theorem support.

\paragraph{Roadmap.}
Section~\ref{sec:primitive-roles} fixes the theorem objects and primitive-role
assignment. Section~\ref{sec:fixed-package} proves the base and bridge theorems.
Section~\ref{sec:frozen-frontier} proves the comparative transfer and
restricted-alignment results on frozen slices. Section~\ref{sec:conditional-lift}
states the main conditional arithmetic theorem and its assumption box.
Section~\ref{sec:vendor-support} explains the supporting-only role of the
vendor-backed evidence. Section~\ref{sec:scope-limits} records the manuscript's
perimeter, including continuity with earlier Six-Birds work and the claims that
remain out of scope.

The theorem-bearing center of the manuscript is the conditional arithmetic lift of Section~\ref{sec:conditional-lift}; the earlier sections build the internal ladder on which it depends, and the later sections delimit its empirical and conceptual perimeter.

\begin{table}[t]
\caption{The theorem ladder used in the paper.}
\label{tab:theorem-ladder}
\centering
\small
\begin{tabularx}{\textwidth}{@{}l X X l@{}}
\toprule
Level & Structural claim & Lean theorem surface & Status \\
\midrule
Base & Fixed-package saturation and no persistent growth under self-iteration & \lid{closure\_iterate\_saturates\_one\_step};\newline \lid{frozen\_package\_no\_persistent\_growth} & in-house \\[6pt]
Bridge & Any genuine strict-growth witness forces package change & \lid{persistent\_growth\_witness\_implies\_package\_change};\newline \lid{bridge\_candidate\_holds\_in\_frozen\_idempotent\_setting} & in-house \\[6pt]
Comparative & Efficiency transfers across cone-equivalent representatives on a frozen slice; restricted alignment holds on the support cone & \lid{frontier\_efficiency\_transfer\_to\_cone\_equivalent};\newline \lid{restricted\_in\_house\_alignment\_lemma} & in-house \\[6pt]
Conditional & Frontier-efficient arithmetic growth admits a canonical representative on the relevant cone under the explicit external contract & \lid{paper\_main\_conditional\_arithmetic\_frontier\_canonicality} & conditional \\
\bottomrule
\end{tabularx}
\end{table}
